\Needspace{4\baselineskip}
\begin{Example}{The heat equation}
\label{ex:heat-fundamental-solution}
For equations which are naturally posed as initial-value problems, a different definition of fundamental solution is used. Consider the equation
\begin{equation}\label{eq:first-order-time-evolution-fundamental-equation}
u_t(\bm{x},t)=L(D)u(\bm{x},t),\quad\bm{x}\in\mathbb{R}^m,\ t>0,
\end{equation}
where $L$ is a constant coefficient differential operator on $\mathbb{R}^m$. Instead of regarding $u$ as a distribution on $\mathbb{R}^m\times(0,\infty)$, we shall in the following regard $u$ as a distribution on $\mathbb{R}^m$, depending on $t$ as a parameter. We say that $u$ depends continuously on $t$ if
\begin{equation}\label{eq:time-parameter-distribution-test-pairing}
\int_{\mathbb{R}^m}u(\bm{x},t)\phi(\bm{x})\,\mathrm{d}\bm{x}
\end{equation}
is continuous in $t$ for every $\phi\in\mathcal{D}(\mathbb{R}^m)$ and we say that $u$ is differentiable with respect to $t$ if \eqref{eq:time-parameter-distribution-test-pairing} is differentiable for every $\phi$. If $u$ is differentiable with respect to $t$, the derivative is also a distribution on $\mathbb{R}^m$; this follows from the representation of the derivative as a limit of difference quotients and the completeness of the space of distributions.
\end{Example}

If $u(\bm{x},t)$ is a distribution on $\mathbb{R}^m$ depending continuously on $t>0$, we can also regard $u$ as a distribution on $\mathbb{R}^m\times(0,\infty)$. This is because every test function $\phi(\bm{x},t)\in\mathcal{D}(\mathbb{R}^m\times(0,\infty))$ can also be thought of as a test function on $\mathbb{R}^m$ which depends continuously on the parameter $t$. Because of Lemma~\ref{lem:distribution-test-sequence-pairing}, this makes
\begin{equation}\label{eq:time-dependent-test-distribution-pairing}
\int_{\mathbb{R}^m}u(\bm{x},t)\phi(\bm{x},t)\,\mathrm{d}\bm{x}
\end{equation}
a continuous function of $t$; hence
\begin{equation}\label{eq:time-dependent-distribution-spacetime-pairing}
\int_0^\infty\int_{\mathbb{R}^m}u(\bm{x},t)\phi(\bm{x},t)\,\mathrm{d}\bm{x}\,\mathrm{d}t
\end{equation}
exists. Hence $u$ defines a linear functional on $\mathcal{D}(\mathbb{R}^m\times(0,\infty))$; the continuity of this functional can be deduced, for example, by representing the outer integral in \eqref{eq:time-dependent-distribution-spacetime-pairing} as a limit of Riemann sums and using the completeness of the space of distributions.

We are now ready to define a fundamental solution.

\begin{Definition}
\label{def:evolution-fundamental-solution}
We call $G:[0,\infty)\to\mathcal{D}^{\prime}(\mathbb{R}^m)$ a \textbf{fundamental solution} of \eqref{eq:first-order-time-evolution-fundamental-equation} if $G$ is continuously differentiable on $[0,\infty)$, and moreover, $G$ satisfies \eqref{eq:first-order-time-evolution-fundamental-equation} with the initial condition $G(\bm{x},0)=\delta(\bm{x})$.
\end{Definition}

In this definition, we think of $u_t$ in \eqref{eq:first-order-time-evolution-fundamental-equation} as differentiation with respect to the parameter $t$. Nevertheless, it is easy to show that $G$ also satisfies \eqref{eq:first-order-time-evolution-fundamental-equation} in the sense of distributions on $\mathbb{R}^m\times(0,\infty)$.

A solution of the inhomogeneous problem
\begin{equation}\label{eq:first-order-evolution-inhomogeneous-initial-problem}
u_t=L(D)u+f(\bm{x},t),\quad u(\bm{x},0)=u_0(\bm{x}),
\end{equation}
where $f$ and $u_0$ have compact support, and $f$ is continuous from $[0,\infty)$ to $\mathcal{D}^{\prime}(\mathbb{R}^m)$, can now be represented as follows:
\begin{equation}\label{eq:first-order-evolution-fundamental-solution-formula}
\begin{aligned}
u(\bm{x},t)&=\int_{\mathbb{R}^m}G(\bm{x}-\bm{y},t)u_0(\bm{y})\,\mathrm{d}\bm{y}\\
&\quad+\int_0^t\int_{\mathbb{R}^m}G(\bm{x}-\bm{y},t-s)f(\bm{y},s)\,\mathrm{d}\bm{y}\,\mathrm{d}s.
\end{aligned}
\end{equation}
The reader should verify that this is indeed a solution (cf. Problem~\ref{ex:inhomogeneous-fundamental-solution-formula}).

We now consider the heat equation in one space dimension. Problem~\ref{ex:heat-step-cauchy-problem} asks for the solution of the problem
\begin{equation}\label{eq:heat-equation-heaviside-initial-problem}
\begin{aligned}
u_t&=u_{xx},\quad x\in\mathbb{R},\ t>0,\\
u(x,0)&=H(x),\quad x\in\mathbb{R}.
\end{aligned}
\end{equation}
The solution can be found by the ansatz $u(x,t)=\phi(x/\sqrt t)$, which reduces the problem to an ODE. The result is
\begin{equation}\label{eq:heat-heaviside-initial-solution}
u(x,t)=\dfrac1{2\sqrt\pi}\int_{-\infty}^{x/\sqrt t}\exp\left(-\dfrac{v^2}4\right)\,\mathrm{d}v,
\end{equation}
To obtain the fundamental solution, we simply need to differentiate with respect to $x$. We thus obtain
\begin{equation}\label{eq:heat-fundamental-solution-one-dimension}
G(x,t)=\dfrac1{2\sqrt{\pi t}}\exp\left(-\dfrac{x^2}{4t}\right).
\end{equation}
The fundamental solution for the heat equation in several dimensions can be obtained as a direct product:
\begin{equation}\label{eq:heat-fundamental-solution-multiple-dimensions}
G(\bm{x},t)=\left(\dfrac1{2\sqrt{\pi t}}\right)^m\exp\left(-\dfrac{|\bm{x}|^2}{4t}\right).
\end{equation}

